\documentclass{article}
\usepackage[UTF8]{ctex}
\usepackage{amsmath,amssymb,amsfonts}
\usepackage{enumitem}
\usepackage{geometry}
\geometry{a4paper, margin=1in}
% === 错题本排版优化（不参与题目标识） ===
\setlength{\parindent}{0pt}
\setlength{\parskip}{0.35em plus 0.10em minus 0.08em}
\setlength{\emergencystretch}{3em}
\setlist{itemsep=0.25em, topsep=0.40em, parsep=0pt}
\allowdisplaybreaks[2]
% === 错题本排版优化结束 ===
\begin{document}
\section*{西安电子科技大学2026年高等代数真题}
\section*{182．西安电子科技大学2026年高等代数试卷}

\subsection{题目1}
1．（10 分）设 \(a_{1}, a_{2}, \cdots, a_{n}\) 是两两不同的非零实数，计算行列式

\[
\left|\begin{array}{cccc}
1 & 1 & \cdots & 1 \\
a_{1} & a_{2} & \cdots & a_{n} \\
a_{1}^{2} & a_{2}^{2} & \cdots & a_{n}^{2} \\
\vdots & \vdots & & \vdots \\
a_{1}^{n-2} & a_{2}^{n-2} & \cdots & a_{n}^{n-2} \\
a_{2} a_{3} \cdots a_{n} & a_{1} a_{3} \cdots a_{n} & \cdots & a_{1} a_{2} \cdots a_{n-1}
\end{array}\right|
\]

\subsection{题目2}
2．（15 分）已知 \(n(n \geq 2)\) 个实系数多项式 \(f_{i}(x)\) 的次数 \(d_{i}(i=1,2, \cdots, n)\) 满足 \(\sum_{i=1}^{n} d_{i}<\frac{n(n-1)}{2}\) ，证明： \(f_{1}(x), f_{2}(x), \cdots, f_{n}(x)\) 在 \(\mathbb{R}\) 上线性无关．

\subsection{题目3}
3．（15 分）设 \(A=\left(\alpha_{1}, \alpha_{2}, \alpha_{3}, \alpha_{4}\right)\) 为 4 阶矩阵，\(\beta\) 为 4 维列向量，方程组 \(A X=\beta\) 的通解为

\[
(1,2,2,1)^{\mathrm{T}}+k(1,-2,4,0)^{\mathrm{T}} .
\]

记 \(B=\left(\alpha_{3}, \alpha_{2}, \alpha_{1}, \beta-\alpha_{4}\right)\) ，求方程组 \(B X=\alpha_{1}-\alpha_{2}\) 的通解。

\subsection{题目4}
4．（15 分）设 \(A\) 为 \(n\) 阶矩阵，满足 \(A^{3}-2 A^{2}-A+2 E=O\) ，证明：

\[
r(A-E)+r(A+E)+r(A-2 E)=2 n .
\]

\subsection{题目5}
5．（10 分）设 \(n\) 阶复矩阵矩阵 \(A\) 的特征多项式为 \(f(\lambda)\) ，若（ \(\left.f(\lambda), f^{\prime}(\lambda)\right)=\lambda\) ，且 \(r(A)=r\left(A^{2}\right)\) ，证明：\(A\) 可相似对角化。

\subsection{题目6}
6．（10 分）已知 \(f\left(x_{1}, x_{2}, \cdots, x_{n}\right)\) 是一个秩为 \(n\) 符号差为 \(s\) 的实二次型，证明：存在维数为 \(\frac{n-|s|}{2}\) 的 \(\mathbb{R}^{n}\)的子空间 \(V_{1}\) ，使得对任意的 \(X=\left(x_{1}, x_{2}, \cdots, x_{n}\right)^{\mathrm{T}} \in V_{1}\) ，有 \(f(X)=0\) ．

\subsection{题目7}
7．（15 分）已知 \(A\) 为实对称矩阵，\(B\) 为正定矩阵，\(A-B\) 为半正定矩阵。证明：

\begin{itemize}
  \item[（1）] \(f(\lambda)=|A-\lambda B|=0\) 的根大于等于1。
  \item[（2）] \(|A| \geq|B|\) ．
\end{itemize}

\subsection{题目8}
8．（10 分）设 \(A\) 为数域 \(P\) 上的 \(n\) 阶矩阵，\(r(A)=n-1, A^{*}\) 为 \(A\) 的伴随矩阵，记 \(A X=0\) 与 \(A^{*} X=0\)的解空间分别为 \(V_{A}, V_{A^{*}}\) ，证明：\(P^{n}=V_{A} \oplus V_{A^{*}}\) 的充要条件是 \(\operatorname{tr}\left(A^{*}\right) \neq 0\) 。

\subsection{题目9}
9．（15 分）记 \(n\) 阶矩阵

\[
A=\left(\begin{array}{ccccc}
2 & 1 & & & \\
1 & 4 & 1 & & \\
& \ddots & \ddots & \ddots & \\
& & 1 & 2(n-1) & 1 \\
& & & 1 & 2 n
\end{array}\right) .
\]

\begin{itemize}
  \item[（1）] 对任意 \(n\) 阶矩阵 \(B\) ，若 \(A B=B A\) ，证明：\(B\) 可对角化．
  \item[（2）] \(C(A)=\{B \mid A B=B A\}\) ，求 \(C(A)\) 的维数．
\end{itemize}

\subsection{题目10}
10．（15 分）设 \(V=\left\{a x^{2}+b x+c \mid a, b, c \in \mathbb{R}\right\}\) ，定义 \((f(x), g(x))=\int_{0}^{1} f(x) g(x) \mathrm{d} x\) ．

\begin{itemize}
  \item[（1）] 证明：\(V\) 上的线性变换 \(\sigma: f(x) \mapsto f(1-x)\) 是对称变换．
  \item[（2）] 求 \(\sigma\) 在基 \(1, x, x^{2}\) 下的矩阵．
\end{itemize}

\subsection{题目11}
11．（20 分）\(\sigma, \tau\) 为 \(n\) 维欧氏空间 \(V\) 上的线性变换，对任意的 \(\alpha, \beta \in V\) ，若

\[
(\sigma(\alpha), \beta)=(\alpha, \tau(\beta)) .
\]

则称 \(\tau\) 与 \(\sigma\) 互为共轭变换．

\begin{itemize}
  \item[（1）] 证明：\(\sigma\) 与 \(\tau\) 互为共轭变换的充要条件是 \(\sigma, \tau\) 在标准正交基下的矩阵互为转置矩阵。
  \item[（2）] 若 \(\sigma, \tau\) 互为共轭变换，证明：\(\sigma\) 的值域是 \(\tau\) 核空间的正交补，即 \(\operatorname{Im} \sigma=(\operatorname{Ker} \tau)^{\perp}\) 。
\end{itemize}
\end{document}
